% TODO make hw stylesheet
\documentclass[10pt]{article}
\usepackage[letterpaper,top=1in,left=1in,right=1in]{geometry}
\usepackage[dvipsnames,svgnames]{xcolor}
\usepackage[shortlabels]{enumitem}
\usepackage[framemethod=TikZ]{mdframed}
\usepackage{amsmath,amssymb,amsthm}
\usepackage[colorlinks]{hyperref}
\usepackage{multicol}
\usepackage{tabularx}
\usepackage{bbm}

\hypersetup{
    colorlinks=true,
    linkcolor=blue,
    filecolor=magenta,
    urlcolor=magenta,
    pdftitle={MAT 344 HW}
}

\usepackage{mathtools}
\usepackage{thmtools}
\usepackage[textsize=footnotesize]{todonotes}
\usepackage{titling}
\usepackage{cancel}

\reversemarginpar
\mdfdefinestyle{mdbluebox}{roundcorner=10pt,innerbottommargin=9pt,
linecolor=blue,backgroundcolor=TealBlue!5,}
\declaretheoremstyle[headfont=\sffamily\bfseries\color{MidnightBlue},
mdframed={style=mdbluebox},]{thmbluebox}

\mdfdefinestyle{mdbloobox}{frametitlefont=\bfseries,innerbottommargin=8pt,
nobreak=true,backgroundcolor=TealBlue!5,linecolor=blue,}
\declaretheoremstyle[headfont=\bfseries\color{MidnightBlue},
mdframed={style=mdbloobox},headpunct={\\[3pt]},postheadspace=0pt,]{thmbloobox}

\mdfdefinestyle{mdredbox}{frametitlefont=\bfseries,innerbottommargin=8pt,
nobreak=true,backgroundcolor=Salmon!5,linecolor=RawSienna,}
\declaretheoremstyle[headfont=\bfseries\color{RawSienna},
mdframed={style=mdredbox},headpunct={\\[3pt]},postheadspace=0pt,]{thmredbox}

\mdfdefinestyle{mdgreenbox}{linecolor=ForestGreen,backgroundcolor=ForestGreen!5,
linewidth=2pt,rightline=false,leftline=true,topline=false,bottomline=false,}
\declaretheoremstyle[headfont=\bfseries\sffamily\color{ForestGreen!70!black},
mdframed={style=mdgreenbox},headpunct={ --- },]{thmgreenbox}

\mdfdefinestyle{mdorangebox}{linecolor=RedOrange,backgroundcolor=RedOrange!5,
linewidth=2pt,rightline=false,leftline=true,topline=false,bottomline=false,}
\declaretheoremstyle[headfont=\bfseries\sffamily\color{RedOrange!70!black},
mdframed={style=mdorangebox},headpunct={ --- },]{thmorangebox}

\mdfdefinestyle{mdblackbox}{linecolor=black,backgroundcolor=RedViolet!5!gray!5,
linewidth=3pt,rightline=false,leftline=true,topline=false,bottomline=false,}
\declaretheoremstyle[mdframed={style=mdblackbox}]{thmblackbox}

\declaretheoremstyle[spaceabove=6pt,spacebelow=6pt]{boldhead}
\declaretheoremstyle[spaceabove=6pt,spacebelow=6pt,headfont=\normalfont\itshape]{italichead}

\declaretheorem[style=boldhead,name=Problem]{problem}
\declaretheorem[style=boldhead,name=Problem,numbered=no]{problem*}
\declaretheorem[style=boldhead,name=Setup]{setup} % for config geo
\declaretheorem[style=boldhead,name=Example,sibling=problem]{example}
\declaretheorem[style=boldhead,name=Theorem,sibling=problem]{theorem}
\declaretheorem[style=boldhead,name=Corollary,sibling=theorem]{corollary}
\declaretheorem[style=boldhead,name=Proposition,sibling=theorem]{prop}
\declaretheorem[style=boldhead,name=Lemma,numberwithin=problem]{lemma}
\declaretheorem[style=boldhead,name=Theorem,numbered=no]{theorem*}
\declaretheorem[style=boldhead,name=Lemma,numbered=no]{lemma*}
\declaretheorem[style=boldhead,name=Claim,numberwithin=problem]{claim}
\declaretheorem[style=boldhead,name=Claim,numbered=no]{claim*}
\declaretheorem[style=boldhead,name=Remark,numberwithin=problem]{remark}
\declaretheorem[style=boldhead,name=Remark,numbered=no]{remark*}
\declaretheorem[style=boldhead,name=Definition,sibling=theorem]{definition}
\declaretheorem[style=boldhead,name=Definition,numbered=no]{definition*}

\providecommand{\ol}{\overline}
\providecommand{\ceil}[1]{\left\lceil #1 \right\rceil}
\providecommand{\floor}[1]{\left\lfloor #1 \right\rfloor}
\newcommand{\dd}{\mathrm{d}}
\providecommand{\ul}{\underline}
\providecommand{\eps}{\varepsilon}
\providecommand{\half}{\frac{1}{2}}
\providecommand{\CC}{\mathbb C}
\providecommand{\FF}{\mathbb F}
\providecommand{\II}{\mathbb I}
\providecommand{\NN}{\mathbb N}
\providecommand{\QQ}{\mathbb Q}
\providecommand{\RR}{\mathbb R}
\providecommand{\ZZ}{\mathbb Z}
\providecommand{\ind}{\mathbbm 1}
\providecommand{\dg}{^\circ}
\providecommand{\ii}{\item}
\providecommand{\setdiff}{\smallsetminus}
\providecommand{\alert}[1]{{\sffamily\textbf{\textcolor{blue}{#1}}}}
\providecommand{\inv}{^{-1}}
\providecommand{\mb}[1]{\mathbf{#1}}
\newcommand{\what}{\widehat}

\newenvironment{soln}{\begin{proof}[Solution]}{\end{proof}}

\providecommand{\pow}{\operatorname{Pow}}
\providecommand{\vocab}[1]{{\textbf{\textcolor{ForestGreen}{#1}}}}
\providecommand{\lcm}{\operatorname{lcm}}
\providecommand{\abs}[1]{\left\lvert #1\right\rvert}
\providecommand{\smabs}[1]{\lvert #1\rvert}
\providecommand{\innerprod}[1]{\langle\!\langle #1\rangle\!\rangle}
\providecommand{\norm}[1]{\left\| #1\right\|}
\providecommand{\smnorm}[1]{\| #1\|}
\providecommand{\gr}{\operatorname{gr}}
\providecommand{\im}{\operatorname{im}}
\providecommand{\paren}[1]{\left(#1\right)}

\usepackage{mathrsfs}

% place title at top right
\setlength{\droptitle}{-8ex}
\pretitle{\begin{flushright}\Large\bfseries}
\posttitle{\par\end{flushright}}
\preauthor{\begin{flushright}}
\postauthor{\end{flushright}}
\predate{\begin{flushright}}
\postdate{\end{flushright}}

\title{MAT 344 HW08}
\author{\large Aditya Pahuja}
\date{\large Due 04/14}
\begin{document}
\maketitle

\begin{problem}
    [p.178\#2]
    Show that the series
    \begin{equation*}
	\zeta(z) = \sum_{n=1}^\infty n^{-z}
    \end{equation*}
    converges for $\operatorname{Re} z > 1$,
    and represent its derivative in series form.
\end{problem}

\begin{proof}
    For any $c > 1$, if $\operatorname{Re} z \ge c$, then
    $\smabs{n^{-z}} \le n^{-c}$.
    Since $\sum_{n=1}^\infty n^{-c} < \infty$,
    the Weierstrass $M$-test implies that $\zeta(z)$ converges
    uniformly on $\mathbb H_c = \{z\in\CC : \operatorname{Re} z\ge c\}$.
    Since any compact subset of this domain of convergence
    is contained in $\mathbb H_c$ for some $c$,
    Theorem 1 on p.176 implies that the derivative of $\zeta$
    can be taken termwise, so
    \begin{equation*}
	\boxed{\zeta'(z) = \sum_{n=1}^\infty -n^{-z}\log n.}\qedhere
    \end{equation*}
\end{proof}

\begin{problem}
    [p.190\#1]
    Comparing coefficients in the Laurent developments of $\cot(\pi z)$
    and of its expression as a sum of partial fractions,
    find the values of
    \begin{equation*}
	\sum_{n=1}^\infty \frac 1{n^2},\qquad
	\sum_{n=1}^\infty \frac 1{n^4},\qquad
	\sum_{n=1}^\infty \frac 1{n^6}.
    \end{equation*}
\end{problem}

\begin{soln}
    Equation (10) on p.189 says that
    \begin{equation*}
	\pi\cot(\pi z)
	= \frac 1z + \sum_{n=1}^\infty \frac{2z}{z^2 - n^2}
	= \frac 1z - \sum_{n=1}^\infty
	\frac 1{n^2}\cdot \frac{2z}{1 - (\frac zn)^2}
	= \frac 1z - \sum_{n=1}^\infty \sum_{k=0}^\infty
	\frac{2z}{n^2}\cdot \frac{z^{2k}}{n^{2k}}
	= \frac 1z - \sum_{k=1}^\infty 2\zeta(2k)z^{2k-1}.
    \end{equation*}

    On the other hand, $\pi\cot(\pi z) = \frac{\pi\cos(\pi z)}{\sin(\pi z)}$,
    so the Laurent series for $\pi\cot(\pi z)$ around zero is given by
    \begin{align*}
	\pi\cdot
	\frac{1 - \frac{(\pi z)^2}{2!} + \frac{(\pi z)^4}{4!} - \cdots}
	{\pi z - \frac{(\pi z)^3}{3!} + \frac{(\pi z)^5}{5!} - \cdots}
	&= \frac 1z\cdot
	\frac{1 - \frac{(\pi z)^2}{2!} + \frac{(\pi z)^4}{4!} - \cdots}
	{1 - \frac{(\pi z)^2}{3!} + \frac{(\pi z)^4}{5!} - \cdots}\\
	&= \frac 1z\paren{
	    1 - \frac{(\pi z)^2}{2!} + \frac{(\pi z)^4}{4!} -
	    \frac{(\pi z)^6}{6!}\cdots
	}
	\paren{
	    1 + \frac{\pi^2 z^2}{6} + \frac{7\pi^4 z^4}{360} +
	    \frac{31\pi^6z^6}{15120} + \cdots
	}.
    \end{align*}
    Comparing coefficients between the the two equations tells us that
    \begin{gather*}
	-2\zeta(2) = \frac{-\pi^2}2 + \frac{\pi^2}6
	= -\frac{\pi^2}3\implies = \boxed{\zeta(2) = \frac{\pi^2}6}\\
	-2\zeta(4) = -\frac{\pi^4}{12} + \frac{7\pi^4}{360}
	+ \frac{\pi^4}{24} = -\frac{\pi^4}{45}
	\implies \boxed{\zeta(4) = \frac{\pi^4}{90}}\\
	-2\zeta(6) = \frac{31\pi^6}{15120}
	- \frac{-7}{720} + \frac 1{144} - \frac 1{720}
	= -\frac{2\pi^6}{945}
	\implies \boxed{\zeta(6) = \frac{\pi^6}{945}}\qedhere
    \end{gather*}
\end{soln}

\begin{problem}
    [p.193\#1]
    Show that
    \begin{equation*}
	\prod_{n=2}^\infty \paren{1 - \frac 1{n^2}} = \half.
    \end{equation*}
\end{problem}

\begin{soln}
    This follows from
    \begin{equation*}
	\lim_{N\to\infty}\prod_{n=2}^N \paren{1 - \frac 1{n^2}}
	= \lim_{N\to\infty}\prod_{n=2}^N \frac{(n + 1)(n - 1)}{n^2}
	= \lim_{N\to\infty} \half\cdot \frac{N+1}N = \boxed\half
    \end{equation*}
    where the second-to-last equality is true because for $2 < n < N$,
    the $n^2$ term in the denominator is canceled by
    the two $n$s in the numerators
    of the previous and next terms of the product.
\end{soln}

\begin{problem}
    [p.193\#2]
    Prove that for $\abs z < 1$,
    \begin{equation*}
	(1 + z)(1 + z^2)(1 + z^4)(1 + z^8)\cdots = \frac 1{1-z}.
    \end{equation*}
\end{problem}

\begin{soln}
    Observe that the product of the first $N$ terms on the left-hand side
    are equal to $\frac{1 - z^{2^N}}{1-z}$ by difference of squares;
    clearly this is true for $N = 1$, and
    \begin{equation*}
	\frac{1 - z^{2^N}}{1-z}\cdot (1 + z^{2^N})
	= \frac{1 - z^{2^{N+1}}}{1 - z},
    \end{equation*}
    so the claim is true for all $N$. Then $z^{2^{N+1}}$ goes to zero
    when $N$ grows arbitrarily large because $\abs z < 1$, so
    \begin{equation*}
	\prod_{n=0}^\infty(1 + z^{2^n})
	= \lim_{N\to\infty}\prod_{n=0}^N (1 + z^{2^n})
	= \lim_{N\to\infty} \frac{1 - z^{2^{N+1}}}{1 - z}
	= \boxed{\frac 1{1-z}}.\qedhere
    \end{equation*}
\end{soln}

\begin{problem}
    [p.193\#3]
    Prove that
    \begin{equation*}
	\prod_{n=1}^\infty \paren{1 + \frac zn} e^{-z/n}
    \end{equation*}
    converges absolutely and uniformly on every compact set.
\end{problem}

\begin{soln}
    Expanding $e^{-z/n}$ as a power series and multiplying, we get
    \begin{equation*}
	\prod_{n=1}^\infty \paren{
	    1 + \sum_{k=2}^\infty (-1)^k\paren{\frac 1{(k-1)!} - \frac 1{k!}}
	    \paren{\frac zn}^k
	},
    \end{equation*}
    which converges absolutely and uniformly on a given compact set $K$
    if and only if the sum
    \begin{equation*}
	\sum_{n=1}^\infty
	\paren{\sum_{k=2}^\infty (-1)^k\paren{\frac 1{(k-1)!}
	- \frac 1{k!}} \paren{\frac zn}^k}
    \end{equation*}
    also converges absolutely and uniformly on $K$.

    Since $K$ is compact, it is bounded,
    say $\abs z < N$ for some integer $N$ whenever $z\in K$.
    Discarding finitely many terms in the series under consideration
    does not affect the nature of its convergence, so we consider
    the sum starting at $n = N$ instead. Then,
    we get both absolute and uniform convergence
    by the following chain of inequalities:
    \begin{align*}
	\sum_{n=1}^\infty
	\abs{\sum_{k=2}^\infty (-1)^k\paren{\frac 1{(k-1)!}
	- \frac 1{k!}} \paren{\frac zn}^k}
	&\le \sum_{n=N}^\infty\sum_{k=2}^\infty
	\paren{\frac 1{(k-1)!} - \frac 1{k!}}
	\frac{\abs z^k}{n^k}\\
	&\le \sum_{n=N}^\infty\sum_{k=2}^\infty
	\paren{\frac 1{(k-1)!} - \frac 1{k!}} \frac{N^2}{n^2}
	= \sum_{n=N}^\infty \frac{N^2}{n^2} < \infty.\qedhere
    \end{align*}
\end{soln}

\begin{problem}
    [p.197\#1]
    Suppose that $a_n\to\infty$ and that the $A_n$
    are arbitrary complex numbers. Show that there exists
    an entire function $f(z)$ which satisfies $f(a_n) = A_n$.
\end{problem}

\begin{soln}
    Let $g(z)$ have simple zeroes at each $a_n$ and consider
    \begin{equation*}
	f(z) = \sum_{n=1}^\infty
	g(z)\frac{e^{\gamma_n(z-a_n)}}{z-a_n}\cdot \frac{A_n}{g'(a_n)}.
    \end{equation*}
    for some constants $\gamma_n$ yet to be determined.
    Certainly if the sum converges locally uniformly
    then $\lim_{z\to a_n} f(z)$ is $A_n$, since
    \begin{equation*}
	\lim_{z\to a_n}\frac{g(z)}{(z-a_n)g'(a_n)}\cdot e^{\gamma_n(z-a_n)}
	= \lim_{z\to a_n}\frac{g(z) - g(a_n)}{(z-a_n)g'(a_n)}
	\cdot e^{\gamma_n(z-a_n)}= 1
    \end{equation*}
    while all the other addends go to zero
    because the factors excluding $g(z)$ would be bounded
    (using uniform convergence to swap sum and limit).

    not sure how to pick $\gamma_n$ :(
\end{soln}

\begin{problem}
    [p.198\#3]
    What is the genus of $\cos(\sqrt z)$?
\end{problem}

\begin{soln}
    Since $\cos(z)$ is a translate of $\sin(z)$,
    it has the same genus as $\sin(z)$, so
    \begin{equation*}
	\cos(z) = e^{g(z)}
	\prod_{n\in\ZZ}\paren{1 - \frac{2z}{\pi(2n-1)}}e^{2z/(\pi(2n-1))}.
    \end{equation*}
    The same strategy as the one employed by Ahlfors
    to conclude that the product converges uniformly
    and that $g(z)$ is a constant in the product expansion for
    $\sin(z)$ also can be used here to show that $g(z)$ is constant here.
    Pairing factors of the product with positive and negative coefficients
    shows that
    \begin{equation*}
	\cos(z) = C\cdot\prod_{n=1}^\infty
	\paren{1 - \frac{4z^2}{\pi^2(2n-1)^2}}
    \end{equation*}
    where the $e^\bullet$ terms cancel each other in each pairing.
    This implies that
    \begin{equation*}
	\cos(\sqrt z) = C\cdot \prod_{n=1}^\infty
	\paren{1 - \frac{4z}{\pi^2(2n-1)^2}},
    \end{equation*}
    which has the shape of the canonical product for the sequence
    $\{\frac 4{\pi^2(2n-1)^2}\}_{n\in\NN}$, hence has genus \framebox{0}.
\end{soln}










\end{document}
